\documentclass[11pt,a4paper]{article}
\usepackage[utf8]{inputenc}
\usepackage[T1]{fontenc}
\usepackage[french]{babel}
\usepackage{amsmath,amssymb,amsthm}
\usepackage{booktabs}
\usepackage{graphicx}
\usepackage[margin=2.5cm]{geometry}

\title{Analyse factorielle Fama-French par régression}
\author{Ton Nom}
\date{\today}

\begin{document}
\maketitle

\begin{abstract}
Résumé en quelques lignes : question, données, méthode, principal résultat.
\end{abstract}

\section{Introduction}
Question posée et motivation.

\section{Cadre théorique}
\subsection{Le modèle linéaire}
Pour un actif $i$, on pose
\[
  r_{i,t} - r_{f,t} = \alpha_i + \boldsymbol{\beta}_i^{\top} \mathbf{F}_t + \varepsilon_{i,t}.
\]

\subsection{Estimateur des moindres carrés}
Avec $y \in \mathbb{R}^T$ et $X \in \mathbb{R}^{T \times (k+1)}$ de rang plein,
\[
  \hat{\theta} = \arg\min_{\theta} \lVert y - X\theta \rVert^2 = (X^{\top}X)^{-1}X^{\top}y.
\]
\emph{À compléter : dérivation par les équations normales, interprétation géométrique (projection orthogonale).}

\subsection{Inférence robuste}
Sous hétéroscédasticité et autocorrélation, la variance de $\hat{\theta}$ est estimée par Newey-West :
\[
  \widehat{\mathrm{Var}}(\hat{\theta}) = (X^{\top}X)^{-1}\,\hat{S}\,(X^{\top}X)^{-1}.
\]
\emph{À compléter : définition de $\hat{S}$ avec poids de Bartlett.}

\section{Données}
Description des facteurs, de l'échantillon d'actifs, de la période.

\section{Résultats}
\subsection{Comparaison CAPM / FF3 / FF5}
\subsection{Stabilité des coefficients}

\section{Limites}
Biais de survie, période, coûts de transaction, tests multiples.

\section{Conclusion}

\end{document}
